\section{Introduction}
\label{sec:intro}

Persistent homology records which topological features of a growing simplicial complex survive as the complex grows~\cite{edelsbrunner2002topological,zomorodian2005computing,carlsson2009topology}. For a two-step filtration \(X_1\subseteq X_2\) and a degree \(d\), the basic invariant is the persistent Betti number, the rank of the induced map \(H_d(X_1)\to H_d(X_2)\). Dividing it by \(\beta_d(X_1)\) gives a scale-free and directly interpretable statistic,
\begin{equation}
\label{eq:intro-ratio}
q_d(X_1,X_2)=\frac{\rank[H_d(X_1)\to H_d(X_2)]}{\beta_d(X_1)}\in[0,1],
\end{equation}
the fraction of the degree-\(d\) holes of \(X_1\) that are still present in \(X_2\). Lowe, Kim, Bondesan, and Hayakawa~\cite{lowe2026normalized} recently introduced this quantity under the name \emph{normalized harmonic persistence} and asked for its complexity on clique complexes of weighted graphs, the input model of quantum topological data analysis~\cite{lloyd2016quantum,hayakawa2022quantum,gyurik2022towards,berry2024analyzing,mcardle2026streamlined}. Their motivation is that the known quantum algorithms for Betti and persistent Betti numbers return quantities normalized by the number of simplices~\cite{lloyd2016quantum,hayakawa2022quantum,schmidhuber2023limitations,apers2023simple}, whereas a fraction of holes is the statistic one would actually ask for. Under inverse-polynomial spectral-gap promises on the two endpoint Laplacians they proved that a low-energy relaxation of \(q_d\) is \(\mathsf{DQC}_1\)-hard and in \(\BQP\), conjectured that \(q_d\) itself is hard for a perfect-completeness variant \(\SDQCone\) of \(\mathsf{DQC}_1\) (their Conjecture~1), and reduced this conjecture to a \emph{coefficient-sensitive kernel realization} of history Hamiltonians by clique complexes (their Conjecture~2). They describe the missing ingredient, a kernel-preserving simulation of frustration-free Hamiltonians by combinatorial Laplacians, as an important open question.

\paragraph{Why exact kernels are hard to realize.}
The exact-kernel primitive available in the literature is the clique-gadget construction of Crichigno and Kohler~\cite{crichigno2024clique} and King and Kohler~\cite{king2026gapped}: a graph whose clique complex has one hole per computational basis state, together with one weighted gadget per local projector of a frustration-free Hamiltonian \(H\). The gadgets produce the desired zero modes, and the many-gadget spectral analysis of King and Kohler proves a gap when the simulated \(H\) is positive definite. A persistence reduction needs more. It needs two Hamiltonians \(H_{\rm in}\) and \(H_{\rm out}=H_{\rm in}+H_{\rm rej}\) with known but generally large kernels, the exact multiplicity \(\dim\ker\Delta=\dim\ker H\) at both endpoints, an inverse-polynomial gap above the \emph{complete} kernel at both endpoints, and control of the inclusion-induced map. One must rule out every additional geometric zero mode and every small positive eigenvalue, including for chains supported almost entirely on private gadget vertices, with constants that do not depend on the exponentially large padded register.

\paragraph{Results.}
We carry out this realization for the history Hamiltonians of exact circuits over the gate set \(G_2=\{X,\CX,\CCX,H\otimes H\}\). Following Rudolph~\cite{rudolph2026universal}, \(\BQPone^{G_2}\) is the class of promise problems decided with perfect completeness by uniform exact \(G_2\) circuits; by his exact simulation theorem it contains \(\BQPone^{G}\) for every finite gate set \(G\) with entries in a cyclotomic field \(\mathbb Q(\zeta_{2^k})\), for instance Clifford\(+T\). Our source, however, is more general. An \emph{\(\SF^{G_2}\) verifier} (\cref{def:sf}) is an exact \(G_2\) circuit on clean qubits and \(r\) unconstrained qubits with a subspace \(S_x\) of the unconstrained register that is accepted with certainty, every state orthogonal to \(S_x\) accepted with probability at most \(1/3\), and a promise \(\dim S_x/2^{r}\ge a\) or \(\le b\). This is the perfect-subspace analogue of \(\mathsf{DQC}_1\) with the decision imposed on the fraction of perfectly accepted states, the quantity that Lowe et al.\ identify as the power of \(\SDQCone\) and from which their hardness proofs decide the instance; it contains \(\BQPone^{G_2}\) through a three-bit label trick, and it contains \(\SDQCone^{G_2}\) itself for all threshold pairs with \(b<(3a-1)/2\).

\begin{theorem}[Main theorem; \cref{thm:general-hardness,cor:sdqc,cor:unweighted-general}]
\label{thm:intro-main}
For every \(a-b\ge1/\poly\) there is a deterministic polynomial-time reduction that maps an instance \(x\) of an \(\SF^{G_2}(a,b)\) problem to nested unweighted clique complexes \(\Xin\subseteq\Xout\) on a common vertex set, a degree \(d\), and rationals \(\gamma_{\rm in},\gamma_{\rm out}\ge1/\poly(|x|)\) such that the endpoint Laplacians have no eigenvalue in \((0,\gamma_{\rm in})\) and \((0,\gamma_{\rm out})\) respectively,
\[
\beta_d(\Xin)=2^{r},\qquad
\rank[H_d(\Xin)\to H_d(\Xout)]=\dim S_x,\qquad
q_d(\Xin,\Xout)=\frac{\dim S_x}{2^{r}} .
\]
Consequently normalized harmonic persistence with inverse-polynomial endpoint gaps (Problem~9 of~\cite{lowe2026normalized}) is \(\SF^{G_2}\)-hard for weighted and for unweighted clique complexes; in particular it is \(\BQPone^{G_2}\)-hard already when \(\beta_d(\Xin)=8\) and \(q_d\in\{3/4,1/8\}\), and it is \(\SDQCone^{G_2}\)-hard in the fraction form for all thresholds and in the trace form of \cite[Definition~2]{lowe2026normalized} whenever \(b<(3a-1)/2\), with \(\beta_d(\Xin)=2^{q-1}\) initial holes.
\end{theorem}

The second result concerns the algorithmic side. Lowe et al.\ show that the problem is in \(\BQP\) when the uniform mixture over \(\cH_1=\ker\Delta_{1,d}\) can be prepared and the pair \((P_{\cH_2},\cH_1)\) satisfies a large-overlap condition (their Lemma~7). Both conditions hold on our instances.

\begin{theorem}[Hard and in \(\BQP\); \cref{thm:advantage,cor:unweighted-general}]
\label{thm:intro-advantage}
The reduction of \cref{thm:intro-main} can be arranged so that every nonzero eigenvalue of \(P_{\cH_{\rm out}}P_{\cH_{\rm in}}P_{\cH_{\rm out}}\) is at least \(1-1/\poly\), and so that it outputs a polynomial-size circuit preparing the uniform mixture over \(\cH_{\rm in}\) to trace distance \(1/\poly\). Normalized harmonic persistence therefore remains \(\SF^{G_2}\)-hard when restricted to instances on which it is in \(\BQP\).
\end{theorem}

The third result is a counting statement that uses only the classical part of the palette.

\begin{theorem}[Gapped Betti numbers; \cref{thm:gapped-betti,cor:unweighted-general}]
\label{thm:intro-betti}
Computing \(\beta_d(\Cl(G))\) for a graph \(G\), weighted or unweighted, promised that \(\Delta_{\Cl(G),d}\) has no eigenvalue in \((0,\gamma)\) with \(\gamma\ge1/\poly\), is \(\sharpP\)-hard under parsimonious reductions and \(\sharpP\)-complete under weakly parsimonious reductions; the same holds for persistent Betti numbers of gapped filtrations.
\end{theorem}

Crichigno and Kohler prove \(\sharpP\)-hardness of Betti numbers of clique complexes without a gap promise~\cite{crichigno2024clique}, Schmidhuber and Lloyd in the clique-dense regime, again without a gap~\cite{schmidhuber2023limitations}, and Cade and Crichigno with a gap for general chain complexes with local coboundaries~\cite{cade2024complexity}. \Cref{thm:intro-betti} shows that the spectral gap which makes the simplex-normalized Betti number estimable by the quantum algorithms does not make the exact Betti number of a clique complex any easier.

\paragraph{The engine.}
All three theorems rest on one transfer theorem and one quotient theorem.

\begin{theorem}[Kernel transfer; \cref{thm:transfer}, informal]
\label{thm:intro-transfer}
Fix a finite palette of weighted clique gadgets satisfying finitely many exactly checkable local conditions (\cref{ass:local-certificate}), and let \(p\) be the largest number of private vertices in a simplex of the certified fillings. Attach \(t\) gadgets with private vertex weight \(\lambda\) to the register, simulating \(H=\sum_j\Pi_j\) with \(H\succeq g(P_V-P_{K})\), \(K=\ker H\). Then every unit chain \(x\) of the full clique complex satisfies
\begin{equation}
\label{eq:intro-concentration}
\dist(x,K)^2\le C\Bigl[t\lambda^2+\frac{\langle x,\Delta x\rangle}{g\lambda^{2p}}\Bigr],
\end{equation}
with \(C\) depending only on the palette. Hence, for \(\lambda\le c\min\{t^{-1},\eta/\sqrt t\}\) and \(E=c'\eta^2g\lambda^{2p}\),
\[
\dim\ker\Delta=\dim K,\qquad\spec(\Delta)\cap(0,E)=\varnothing,\qquad
\|P_{\ker\Delta}-P_K\|\le\eta/2,
\]
including the case \(K=0\).
\end{theorem}

\Cref{eq:intro-concentration} concerns arbitrary geometric chains, not only vectors of the simulated register. Its shape matters. The structural error \(t\lambda^2\) does not involve \(g\), so the gadget scale \(\lambda\) is chosen from the number of gadgets alone, independently of the logical gap, and the resulting geometric gap \(E\) is linear in \(g\); a direct use of a many-gadget inequality in which \(\lambda\) is proportional to \(g\) would instead produce a high power of \(g\). The only power of \(\lambda\) beyond \(\lambda^2\) is \(\lambda^{2p}\), which comes from a Schur-complement bound through the certified fillings and is independent of the locality of the gadgets; for our palette \(p=5\), and a conservative choice gives \(E=\Omega(\eta^{12}g/t^{10})\). The projector bound is what yields \cref{thm:intro-advantage}.

\begin{theorem}[Quotient naturality; \cref{thm:quotient}]
\label{thm:intro-quotient}
Under the hypotheses of \cref{thm:intro-transfer}, \(H_d(X_A)\cong V/W_A\) with \(W_A=\sum_{j\in A}\ran\Pi_j=K_A^\perp\), and for \(A\subseteq B\) the inclusion \(X_A\subseteq X_B\) induces the natural quotient map \(V/W_A\twoheadrightarrow V/W_B\). In particular \(\rank[H_d(X_A)\to H_d(X_B)]=\dim K_B\).
\end{theorem}

\Cref{thm:intro-transfer,thm:intro-quotient} give the two properties that Conjecture~2 of Lowe et al.\ requests, kernel fidelity with computable gaps and functoriality along the filtration, for the exact history Hamiltonians treated here; \cref{sec:related} states the precise comparison. Functoriality is obtained differently from the way it is formulated there. We never choose harmonic representatives compatibly across the filtration; exact fillings identify each homology group with a quotient of the fixed register cycle space, the induced map is then forced to be the natural quotient, and the projector bound shows a posteriori that the harmonic spaces are small rotations of the logical kernels.

\paragraph{Technical overview.}
Five ideas carry the proofs.

\emph{Finite certificates and the zero-weight limit.} In normalized coordinates the boundary pair of a gadget with private weight \(\lambda\) is affine in \(\lambda\), \(D(\lambda)=D_0+\lambda D_1\), and \(D_0\) is a fixed rational matrix. The certificate (\cref{ass:local-certificate}) asks for four exact facts about each palette member: the boundary intersection \(B_{d_a}(Y_a)\cap V_a\) equals the range of the intended rank-one projector (exact filling); \(\ker D_0=V_a\oplus Q_a\) for an explicit private subspace \(Q_a\); a projected private differential is coercive on \(Q_a\) at first order in \(\lambda\); and its retained outputs are private. All four are integer or modular linear algebra, and \cref{app:palette} records them for the palette of this paper.

\emph{Harmonic padding.} A certificate speaks about a constant-size active register, while the actual gadget is joined with the unaffected register factors, so the padded complex is exponentially large. The join Laplacian is a tensor sum over bidegrees, and the unaffected factors have harmonics only in top degree. We keep the local private data tensored with these outside harmonics and place every nonharmonic outside mode in a sector with a constant gap (\cref{lem:scaling-padding}). This is the step where the constants become independent of the padded dimension.

\emph{Two energies and the filling depth.} The proof of \cref{eq:intro-concentration} separates leakage out of the register from logical energy. Leakage is controlled by the zero-weight floor, the private coercivity, and an interface bound \(\|[L_1\cdots L_t]\|^2=\|\sum_jL_jL_j^*\|\le Ct\lambda^2\) that needs no orthogonality between the register outputs of different gadgets. Logical energy is controlled through \(\Delta^\uparrow\succeq c\lambda^{2p}H^{\rm emb}\), which follows from the Cauchy--Schwarz inequality applied to the weighted certified fillings (\cref{lem:filling-domination}); the exponent is twice the number of private vertices in the deepest filling simplex.

\emph{From estimate to exact multiplicity and overlap.} Every vector of the spectral subspace of \(\Delta\) below \(E\) projects nontrivially onto \(K\), so that subspace has dimension at most \(\dim K\); exact fillings inject \(V/W_A\) into \(H_d(X_A)\), providing \(\dim K\) zero modes. Equality forces the exact multiplicity and the absence of eigenvalues in \((0,E)\). Applied at zero energy, the same estimate says that every harmonic vector is within \(\eta/2\) of \(K\); equal dimensions then give \(\|P_{\ker\Delta}-P_K\|\le\eta/2\), and since \(K_{\rm out}\subseteq K_{\rm in}\), the nonzero singular values of \(P_{\cH_{\rm out}}P_{\cH_{\rm in}}\) are all at least \(1-\eta\).

\emph{Exact sources.} \Cref{prop:history} gives the two history kernels and gaps for any exact circuit with a perfectly accepted subspace and soundness \(1/3\) on its complement, whatever the clean and unconstrained registers are. Three instances matter: a three-bit label turns any \(\BQPone^{G_2}\) verifier into such a source with compressed acceptance operator \(\operatorname{diag}(1,p_x,p_x,p_x,p_x,p_x,0,0)\); a one-clean-qubit verifier is the \(\SDQCone\) source; and a reversible evaluation of a Boolean formula gives the counting source, with \(S\) spanned by the satisfying assignments. Hayakawa's common-copy blow-up~\cite{hayakawa2026unweighted}, applied with the same labeled copy blocks at both levels, transfers everything to unweighted graphs.

\paragraph{Scope.}
The hardness is gate-dependent, as in the known perfect-completeness hardness results for topological data analysis~\cite{crichigno2024clique,king2026gapped,gyurik2026provable}, and as in the definition of \(\SDQCone\) itself. We do not obtain ordinary \(\BQP\)-hardness, gate-independent \(\BQPone\)-hardness, or \(\SDQCone\)-hardness in the trace form for threshold pairs with \(b\ge(3a-1)/2\); the last restriction is unavoidable for any reduction whose target sees only the fraction of perfectly accepted states (\cref{app:source-boundaries}). Complex phases and general integer states are outside the certified palette. The gap is a complexity promise rather than a practical estimate. \Cref{sec:limitations} discusses these points and the open problems they suggest.

\paragraph{Organization.}
\Cref{sec:prelim} fixes conventions and the problem. \Cref{sec:transfer} states the certificate model and the transfer theorem, \cref{sec:quotient} derives quotient naturality, \cref{sec:circuit} instantiates the history construction with a certified palette and proves the hardness and counting theorems, \cref{sec:advantage} proves the containment conditions, and \cref{sec:unweighting} transfers everything to unweighted graphs. \Cref{sec:related} compares with prior work and \cref{sec:limitations} concludes. The appendices contain the complete transfer proof, the history-gap and source calculations, and the finite palette certificates.
